\chapter{Módulo de Automatización (Auto Mode)}
\label{ch:automode}
% ==============================================================

\section{¿Qué problema resuelve?}

Sin este módulo, procesar un lote nuevo de capturas de una cámara requiere repetir a mano,
imagen a imagen, tres pasos ya vistos en capítulos anteriores: descargar las fotos nuevas de
Obscape, alinearlas contra una referencia (capítulo~\ref{ch:batch}) y analizarlas para extraer
la línea de costa (capítulo~\ref{ch:segmentation}). \emph{Auto Mode} automatiza esas tres cosas
de punta a punta para un rango de fechas elegido por el usuario, sin reimplementar ninguna: las
orquesta, llamando directamente a las mismas funciones que usa el flujo manual.

\begin{infobox}[¿Sobre qué se ejecuta --- una imagen / un lote / toda la cámara?]
Sobre un \textbf{lote acotado por fecha}, elegido explícitamente por el usuario (cámara +
rango \emph{desde}/\emph{hasta}) — nunca se descargan ni procesan imágenes fuera de ese rango,
y nunca se lanza automáticamente sin que el usuario pulse \emph{Iniciar}. Dentro de ese rango,
se procesan \textbf{todas} las imágenes nuevas encontradas (las que ya existen en disco se
saltan). No sustituye a la calibración (capítulo~\ref{ch:calibration}), que sigue siendo un
paso manual previo y obligatorio.
\end{infobox}

\section{Requisito previo: la cámara debe estar ya calibrada}

Auto Mode analiza imágenes con \texttt{analyze\_roi()}, el mismo endpoint que usa el botón
manual \emph{Analizar imagen} — y ese endpoint necesita una homografía activa para la cámara
(capítulo~\ref{ch:calibration}, \S4.2). Por eso, al pulsar \emph{Iniciar} el backend comprueba
primero si existe \texttt{calibration/cam\_\{id\}\_H.npy} o un \texttt{H} guardado en el
perfil de la cámara; si no, rechaza la petición con un mensaje explícito en vez de descargar
imágenes que luego no se podrían georreferenciar:

\begin{codebox}[\_is\_calibrated() --- backend/auto\_mode.py]
\begin{minted}{python}
def _is_calibrated(cam_id: int) -> bool:
    if os.path.exists(f"{CALIBRATION_DIR}/cam_{cam_id}_H.npy"):
        return True
    profile = load_json(f"{CALIBRATION_DIR}/cam_{cam_id}_profile.json")
    return bool(profile and profile.get("H"))
\end{minted}
\end{codebox}

La calibración inicial de una cámara \textbf{no} se automatiza: depende de coordenadas UTM de
varillas topografiadas en campo, que no vienen de la API de imágenes de Obscape. Auto Mode
empieza donde termina la calibración, no la sustituye.

\section{Las cuatro fases}

\begin{figure}[H]
\centering
\begin{tikzpicture}[node distance=0.5cm and 1.4cm]
  \node[block, fill=SkyBlue!15] (desc) {1. Descarga\\{\tiny Obscape API}};
  \node[block, fill=Orange!15, right=of desc] (base) {2. Base\\{\tiny referencia de calibración}};
  \node[block, fill=OliveGreen!15, right=of base] (alin) {3. Alineación\\{\tiny reutiliza cap.~\ref{ch:batch}}};
  \node[block, fill=NavyBlue!15, right=of alin] (anal) {4. Análisis\\{\tiny reutiliza cap.~\ref{ch:segmentation}}};

  \draw[arrow] (desc) -- (base);
  \draw[arrow] (base) -- (alin);
  \draw[arrow] (alin) -- (anal);
\end{tikzpicture}
\caption{Pipeline de Auto Mode --- cada fase reutiliza código ya existente, sin reimplementar nada}
\label{fig:automode-pipeline}
\end{figure}

\begin{enumerate}[leftmargin=1.5em]
  \item \textbf{Descarga}: pide a la API de Obscape los puntos de datos de la cámara en
        \([\text{desde}, \text{hasta}]\) y descarga las imágenes que aún no existen en disco
        (mismo criterio de deduplicación que el descargador manual,
        \texttt{acces\_api/obscape\_api.py}), guardándolas planas en
        \texttt{DATA\_DIR/\{carpeta\_cámara\}/} con el nombrado
        \texttt{<epoch>\_<AAAAMMDD>\_<HHMMSS>\_<serie>.jpg} que el resto del backend ya
        reconoce. Si no hay imágenes nuevas en el rango, el job termina ahí con un mensaje
        claro (nada que hacer) en vez de seguir con pasos vacíos.
  \item \textbf{Selección de base}: TODO el lote se alinea contra la \texttt{reference\_image}
        \textbf{persistente} de la cámara — la misma imagen que fija \texttt{/set-reference} y
        usa por defecto el flujo manual (capítulo~\ref{ch:batch}), \emph{nunca} la imagen más
        antigua de este lote concreto. Es un cambio deliberado de 2026-08-11: si cada lote
        semanal se alineara contra su propia base, cada ejecución quedaría en un marco de
        píxeles distinto al que se usó para calcular la homografía H (que se computa una sola
        vez, sobre \texttt{calibrated\_image}) — la H fija aplicada a un marco de píxeles
        distinto introduce un desplazamiento sistemático que ni el RMSE de calibración (solo se
        calcula una vez) ni la confianza de SAM detectan. Si la cámara no tiene
        \texttt{reference\_image} guardada (calibrada sin pasar nunca por
        \texttt{/set-reference}), cae al comportamiento anterior: la imagen más antigua del
        propio lote, y un lote de una sola imagen no se alinea (nada dentro del lote contra lo
        que alinearla).
  \item \textbf{Alineación}: construye un \texttt{BatchJob} igual que el flujo manual de
        alineación masiva y llama directamente a \texttt{batch\_alignment.\_run\_pipeline()} y,
        si el resultado queda \texttt{ready}, a \texttt{commit\_batch()} — sin pantalla de
        revisión intermedia (\S\ref{sec:automode-commit} explica por qué es seguro comitear sin
        que un humano revise cada imagen).
  \item \textbf{Análisis}: para cada imagen del lote (la alineada si el commit tuvo éxito, la
        original si no — \texttt{\_resolve\_camera\_image\_path()} ya resuelve esa prioridad,
        capítulo~\ref{ch:batch}), llama a \texttt{main.analyze\_roi()}: la misma función que
        ejecuta el botón manual \emph{Analizar imagen}. Segmenta, extrae la línea de costa, la
        proyecta a UTM con la homografía activa de la cámara y guarda el GeoJSON resultante.
\end{enumerate}

\section{Por qué el commit automático es seguro sin revisión humana}
\label{sec:automode-commit}

El módulo de alineación masiva (capítulo~\ref{ch:batch}) se diseñó para uso interactivo: un
humano revisa cada imagen alineada y aprueba o descarta antes de comitear. Auto Mode no tiene
ese paso — comitea en cuanto el job de alineación termina. Esto es seguro porque
\texttt{AlignmentResult} fuerza \texttt{approved=False} automáticamente en cualquier resultado
\texttt{"failed"} (ver \texttt{\_\_post\_init\_\_} en el capítulo~\ref{ch:batch}), y
\texttt{commit\_batch()} solo copia al directorio final las imágenes aprobadas: una alineación
fallida queda excluida del commit por sí sola, sin que nadie tenga que marcarla a mano, y el
análisis posterior usa entonces la imagen \textbf{original sin transformar} para esa captura en
concreto (en vez de una alineación mal hecha).

\section{Modelos de dominio}

\begin{codebox}[Dataclasses principales (backend/auto\_mode.py)]
\begin{minted}{python}
@dataclass
class ImageResult:
    filename:      str
    captured_at:   Optional[str] = None
    is_base:       bool = False
    align_status:  str = ""    # "" | "ok" | "failed" | "reference"
    confidence:    float = 0.0
    dry_area_m2:   float = 0.0
    rejected:      bool = False
    reject_reason: str = ""
    error:         str = ""

@dataclass
class AutoJob:
    job_id:           str
    cam_id:            int
    from_date:        str
    to_date:          str
    status:           str = "pending"  # pending|downloading|aligning|analyzing|done|error
    step:             str = ""         # descripción legible del paso actual
    progress_current: int = 0
    progress_total:   int = 0
    downloaded:       List[str] = field(default_factory=list)
    results:          List[ImageResult] = field(default_factory=list)
    error_msg:        str = ""
\end{minted}
\end{codebox}

\section{API REST del módulo}

\begin{table}[H]
\centering
\caption{Endpoints del router \texttt{/api/automode/}}
\label{tab:api-automode}
\begin{tabular}{C{1.5cm} L{5cm} L{6cm}}
\toprule
\textbf{Método} & \textbf{Ruta} & \textbf{Descripción} \\
\midrule
\texttt{POST} & \texttt{/start}            & Valida cámara calibrada y encola el pipeline completo en background \\
\texttt{GET}  & \texttt{/\{id\}/status}    & Polling de progreso (fase actual + contadores) \\
\texttt{GET}  & \texttt{/\{id\}/results}   & Resultado detallado por imagen, solo cuando \texttt{status} es \texttt{done} o \texttt{error} \\
\bottomrule
\end{tabular}
\end{table}

Igual que el módulo de alineación masiva, el estado de cada job vive \textbf{solo en memoria}
(\texttt{\_jobs}, diccionario global) — no persiste entre reinicios del backend, y el frontend
hace polling a \texttt{/status} cada 1,5~s mientras el job está en curso.

\section{Nota de empaquetado: import dinámico y PyInstaller}

\texttt{\_download\_new\_images()} importa \texttt{obscape\_api} dentro de la propia función,
no como import de módulo al principio del fichero (el mismo patrón que usan
\texttt{segmentation\_sam} y compañía en \texttt{main.py}, capítulo~\ref{ch:architecture}). El
análisis estático de PyInstaller no sigue ese tipo de import, así que sin declararlo
explícitamente en \texttt{hiddenimports} (\texttt{backend/desktop\_launcher.spec}) Auto Mode
habría funcionado perfectamente en desarrollo y fallado con \texttt{ModuleNotFoundError} en el
\texttt{.exe} empaquetado — un fallo que solo se ve al probar el instalador final, no en
desarrollo. Corregido añadiendo \texttt{"obscape\_api"} a \texttt{hiddenimports} y
\texttt{acces\_api/} a \texttt{pathex} en el \texttt{.spec}.

% ==============================================================
